\begin{proposition}[The hypotheses satisfied by the compact example]\hypertarget{tag-01BM}{}\label{cdp:st:B3}
\StatementRecord{cdp:st:B3}{proved}{01BM}{\SeriesRef{cdp:st:B2}; \cite{CDP20}}

The threefold $Y$ of Statement \ref{cdp:st:B2} has $H^1(Y,\Z)=0$, $H^2(Y,\Z)\simeq\Z^{11}$, algebraic dimension $3$, and $e(Y)=24$. Its map $g$ is holomorphic, proper, surjective, and has connected fibres. Thus it tests the fibrewise inference in Section 7 of \cite{CDP20}; it does not satisfy the vanishing of $H^2(Y,\Z)$ or algebraic dimension one required by that paper's Theorem 2.2.
\end{proposition}


\begin{proof}
\begin{enumerate}[label={$\langle1\rangle$\arabic*.},leftmargin=*]
\item 
\noindent\textbf{Step 1.} The residue calculation above gives $K_E=h^*\OO(-1)$. Relative duality gives $R^1h_*\OO_E=\OO(-1)$, since every fibre is an integral genus-one curve, and $h_*\OO_E=\OO$. Thus $q(E)=p_g(E)=0$ and $H^0(E,K_E^{\otimes m})=0$ for every positive $m$. Castelnuovo's rationality criterion, in the form $q=P_2=0$ for a smooth projective surface, makes $E$ rational. A rational smooth projective surface is obtained from a minimal rational surface by point blowups; its first integral homology vanishes and its second integral cohomology is free. Euler additivity gives $e(E)=12$, since the twelve nodal fibres each have Euler characteristic one and the smooth elliptic fibre has Euler characteristic zero. Poincare duality gives $b_2(E)=10$. The integral Kunneth isomorphism for $Y=E\times\PP^1$ gives the asserted groups and Euler number. Projectivity and dimension give $a(Y)=3$. Connectedness and surjectivity follow directly from the integral cubic fibres of $h$ and the second product factor.

\end{enumerate}
\end{proof}